\section[Betti 与 de Rham]{Betti 与 de Rham（3 月 25 日）}
\srcpage{43}

在推广到非阿贝尔群之前，先讨论这个对偶故事中 Betti 版本与 de Rham 版本的差别，
这或许有助于澄清一些事情。回忆我们的范式：A 侧关乎拓扑，B 侧关乎代数几何
（域构成一个代数簇）。在 A 侧，我们常常取像局部常值函数这样的东西。
\srcpage{44}

把「局部常值函数」这个想法形式化，至少有三种方式，它们导致对偶范式的不同实现。

（1）第一种是真正的局部常值函数，即沿道路常值的函数。过渡到导出情形，就得到空间的
奇异上同调。这是「Betti 情形」。

（2）「de Rham 情形」假定空间是光滑的，并说局部常值意味着导数为零。过渡到导出情形，
就得到 de Rham 上同调。可以把 $X$ 上局部常值函数这一概念，等同于空间 $X_{\mathrm{dR}}$
上的函数（回忆 $X_{\mathrm{dR}}$ 不是空间，而是某种叠，其中把无穷小邻近的点等同起来）。
当然，在合理的空间上，奇异上同调与 de Rham 上同调是同构的。

（3）关于这个例子我们不打算多讲，但在辛拓扑中，我们把 $X$ 的拓扑与 $T^*X$ 的辛拓扑
联系起来，具体地说就是 Floer 理论。字母 A 的使用往往与辛拓扑相关（例如 A-模型），
尽管这个对偶也有 Betti 与 de Rham 两种进路。本课只关心余切丛。

范畴化时这会更微妙（话说回来，什么不是呢？）。在余维 $2$，我们要指派一个范畴，
它像是「域的空间」上的层范畴。在 B 侧，域是一个代数簇，该取哪种层没有疑问：
拟凝聚层。\footnote{课程后面会推广到凝聚层之外（ind-coherent sheaves）$\QCoh^{!}(\cF)$，
它更擅长处理奇点。若域的空间是光滑的，则两者没有差别。}
在 A 侧，上面那些不同的函数概念提升为不同的层概念。

对 Betti 与 de Rham 两种情形，都会有一个共同的核心，因为这两种进路在某种意义上都是拓扑的。

（1）在 Betti 情形，核心是局部常值有限秩层，即局部同构于 $U\times\C^r$ 的层。
我们沿道路有平行移动，且在 homotopy 意义下良定义。这个范畴记作
$\mathrm{Loc}^{\mathrm{fd}}(X)$。

（2）在 de Rham 情形，我们考虑平坦向量丛。它们同样是局部平凡的，
这是「导数为零」的一种说法。这是对无穷小邻近对象的一种等变性（TODO：？？）。
这个范畴记作 $\mathrm{Conn}^{\flat}(X)$。

（3）在辛情形，你得到某种 Fukaya 范畴，尽管精确化它的方式不止一种。
这个情形我们不打算太操心。\yizhu{原文把这一条（3）印在（1）（2）之前，
编号顺序有误；今按编号移到这里，内容未改。}{A（原文 p.44）}

这两类层是等价的，这是 Riemann–Hilbert 对应的一个婴儿情形。
但它们坐落在相当不同的、更大的层范畴里。

（1）在 Betti 侧有几种想法，但可以取 $\mathrm{Loc}(X)$，即 $X$ 上所有局部常值层构成的范畴。
在阿贝尔范畴层面，它等价于 $X$ 的基本群胚的表示范畴；我们是在道路上指定局部常值条件。
\footnote{后面还会考虑更大的空间范畴：可构造层或反常层，它们对奇点有用。}

（2）在 de Rham 侧，考虑 $D_X$-模构成的范畴，即带平坦联络的拟凝聚层。
它等价于 $\QCoh(X_{\mathrm{dR}})$。

这两个范畴相当不同。用个简单的例子来看，$X=\C^\times$（拓扑上可以想成圆周）。

（1）拓扑上，$\C^\times$ 上的局部常值层等价于 $S^1$ 上的局部常值层，等价于
$B\Z$ 上的局部常值层，等价于 $\Z$ 的表示。也就是说，要了解 $\C^\times$ 上的局部常值层，
只需知道绕过那个非平凡环路时发生了什么。

这等价于 $\mathrm{Mod}_{\C[\Z]}=\mathrm{Mod}_{\C[z,z^{-1}]}=\QCoh(\G_m)$。
这个 $\G_m$ 在 Cartier 对偶意义下对偶于我们出发时的 $\C^\times$；
若从环面出发，就会得到对偶环面上的拟凝聚层。

一些局部常值层的例子是特征标层 $L_\mu$，即单值为 $\mu$ 的秩一局部系统。
在 Cartier 对偶下，它们与 $\G_m$ 上的摩天楼层互换。
\srcpage{45}

（2）de Rham 的故事是 Mellin 变换的一个版本，相当不同。Mellin 变换是群 $\R_+$ 的
Fourier 变换——当然，这个群经指数映射同构于 $\R$，所以它与通常的 Fourier 变换是同一回事，
只是用乘法记号写。若 $f$ 是 $\R_+$ 上的函数，则
\[
\hat{f}(s)=\int_0^\infty f(z)\,z^s\,\frac{dz}{z}.
\]
这里 $z^s$（$s\in\C$）是 $\R_+$ 的特征标；回到 $\R$ 上，它们变成 $e^{zs}$。

$\C^\times$ 上的微分算子环有一个导子 $\partial:=t\partial_t$，其中 $t$ 是坐标：
\[
D_{\C^\times}=\C\langle t,t^{-1},\partial\rangle\big/
\bigl(\partial t=t\partial+t\bigr).
\]
可以交换这些生成元的次序，写成
$D_{\C^\times}\cong\C\langle s,\sigma,\sigma^{-1}\rangle/(s\sigma=\sigma s+\sigma)$。
这看起来像 $\C$ 上的函数，连同按平移作用的移位算子 $\sigma$！也就是说，$\Z$ 通过平移
作用在 $\C$ 上；这些有限差分算子与 $D_{\C^\times}$ 具有相同的关系。
记 $\Delta_{\C}$ 为这个微分算子环。

过渡到模时，$D_{\C^\times}$-模等价于 $\Delta_{\C}$-模，也就是 $\C$ 上带移位作用的
拟凝聚层，或等价地 $\QCoh(\C/\Z)$。

作为例子，有模 $M_\lambda$，它对应（TODO：我漏了什么）$\C^\times$ 上平凡线丛带平坦联络
\[
\nabla=d-\lambda\,\frac{dt}{t}.
\]
于是 $\partial$ 作用为 $\lambda$；$\partial$ 是分次算子，所以这个 ODE 的任何解 $f$
都是 $\lambda$ 次齐次的，例如 $z^\lambda$。

现在对 $n\in\Z$，$M_\lambda\cong M_{\lambda+n}$，因为 $M_\lambda$ 只依赖其单值，即 $\exp(\lambda)$。

TODO：关于 $\C^\times$ 上的 D-模我漏讲了一些东西，它们有形形色色的对象，如 $\delta$-函数。
Betti 与 de Rham 的差别，对应 $\C/\Z$ 与 $\C^\times$ 的差别：二者解析上相同，
但在代数簇上结构相当不同。这是 B 侧两个模型（Betti 与 de Rham）之间的差别；
在 A 侧，它们对应对层的两种不同概念。

回到类域论。我们说过，$C$ 上的阿贝尔局部系统与 $\mathrm{Jac}(C)$ 上的有对应；
且 $C$ 上秩一局部系统等价于 $\Pic(C)$ 上的特征标局部系统。看范畴时，它们显得相当不同。

在 A 侧，空间是 $\Pic(C)=\Z\times\mathrm{Jac}(C)$（还有一个 $B\G_m$ 因子，
在过渡到导出情形之前它无关紧要）。取层（线性化）有两种方式：Betti 版本取局部系统，
de Rham 版本取 D-模。二者都含有共同的核心，但实际上不同：
$\mathrm{Loc}(\Pic(C))$ 等价于 $\pi_1(\mathrm{Jac}(C))=H_1(C;\Z)\cong\Z^{2g}$ 的
$\Z$-分次表示构成的范畴，其中 $g$ 是 $C$ 的亏格。这既包含有限维单值表示，
也包含自由的那些（TODO：这些不也是有限秩吗？），例如万有覆叠表示。

在 de Rham 版本中我们有 D-模，但因为 Jacobian 不是仿射簇，描述起来更烦人：
它不太是模范畴。我们改为取 $\QCoh(\Pic(C))$ 中的对象，连同向量场的作用。
这与 Betti 故事不同，但同样有共同的对象。

而我们得到的类域论在这两个视角下不同：两种情形都有 Fourier–Mukai 对偶。
（TODO：漏了一个，抱歉。总体而言，今天我对于对象何时该是拓扑空间、何时该是代数簇
一直很困惑，也就是 $\C^\times$ 与 $\G_m$ 或 $\C/\Z$ 之类何时有区别。
这可能让笔记相当不清楚。抱歉。）

也就是说，在 Betti 情形下，特征标层 $\chi_L\in\mathrm{Loc}(\Pic(C))$
在 Fourier–Mukai 变换下与 $\QCoh(\mathrm{Loc}_1(C))$ 中 $L$ 处的摩天楼层
$\cO_L$ 互换。

以上是 Betti 侧——在 de Rham 侧，我们从 Mukai 所做的事情开始：导出范畴的等价
\[
\bigl(\QCoh(\mathrm{Jac}(C)),*\bigr)
\xleftrightarrow{\ \simeq\ }
\bigl(\QCoh(\mathrm{Jac}(C)),\otimes\bigr).
\]
我们想把这些范畴放大：例如我们有 D-模。它们是关于 $\mathrm{Jac}(C)$ 的形式群等变的
拟凝聚层；等价地说，有向量场代数作用在 $\mathrm{Jac}(C)$ 上。
\srcpage{46}

这里藏着一些单子性（monadicity）/Barr–Beck 定理，与这个额外的代数结构有关。

另一侧，我们放大到 $\QCoh(\mathrm{Conn}^{\flat}_1(C))$，即 $\mathrm{Jac}(C)$ 上的
拟凝聚层，连同纤维上函数的作用。可以验证 Fourier–Mukai 变换实际上互换这两类层。

有忘却映射 $\mathrm{Conn}^{\flat}_1(C)\to\mathrm{Jac}(C)$，它忘掉联络，其核为
\[
\Gamma\bigl(\mathrm{Jac}(C),\Omega^1\bigr)
\cong T^*_e\mathrm{Jac}(C)\cong H^0(C;\Omega^1),
\]
即忘却映射是一个仿射丛。

Betti 故事比 de Rham 故事初等得多：你可以手算验证等等，但你失去了东西：
在 de Rham 故事中，两边都生活在代数几何里。这意味着 de Rham 理论有一个额外的对称性，
而它在 Betti 故事中被破坏了。

如何增强前式？可以从 $\mathrm{Jac}(C)$ 过渡到其余切丛。因为 $\mathrm{Jac}(C)$ 是群，
其余切丛平凡，并且
\[
T^*\mathrm{Jac}(C)\cong \mathrm{Jac}(C)\times H^0(C;\Omega).
\]
纤维 $H^0(C;\Omega)$ 若想说得花哨些，可以叫秩一 Hitchin 空间：
这是一般而有趣的对象的一个特殊阿贝尔情形。从 $T^*\mathrm{Jac}(C)$ 与
$\mathrm{Jac}(C)\times H^0(C;\Omega)$ 两个视角，都有到 $H^0(C;\Omega)$ 的投影映射，
它们的纤维等价；而这是 Hitchin 可积系统自对偶性的一个特殊情形。

现在对上述等价作形变量子化。有从 $\cO(T^*X)$ 到微分算子
$D_X=\C\langle x_i,\partial_i\rangle/\sim$ 的自然形变。可以这样总结：加一个参数 $\hbar$，
考虑
\[
\C\langle x_i,\partial_i\rangle\big/
\bigl(\partial_i x_j-x_j\partial_i=\delta_{ij}\hbar\bigr).
\]
当 $\hbar=1$ 时这是 $D_X$，当 $\hbar=0$ 时这是 $\cO(T^*X)$；
D-模是余切丛上函数的形变量子化，正是这个意思。

在 B 侧，形变量子化把 $\QCoh(T^*\mathrm{Jac}(C))$ 形变为
$\QCoh(\mathrm{Conn}^{\flat}_1(C))$。

非阿贝尔情形不那么显式，这使事情更难；推广时把阿贝尔情形当作脚手架或直觉，
或许会有帮助。
